% !TeX program = xelatex
% !TeX encoding = UTF-8
\documentclass[
  openany,           % 章节可以从任意页开始
  12pt,              % 正文字号
  a4paper            % A4纸张
]{ctexbook}

% ==================== 宏包引入 ====================
\usepackage{geometry}
\usepackage{fancyhdr}
\usepackage{xeCJK}
\usepackage{booktabs}
\usepackage{longtable}
\usepackage{tabularx}
\usepackage{setspace}
\usepackage{indentfirst}
\usepackage{footmisc}
\usepackage{hyperref}
\usepackage{titlesec}
\usepackage{titletoc}
\usepackage{ctex}
\usepackage{graphicx}
\usepackage{color}
\usepackage{xcolor}
\usepackage{ulem}

% ==================== 页面设置 ====================
\geometry{
  paper=a4paper,
  top=25mm,            % 天头
  bottom=20mm,         % 地脚
  left=25mm,           % 订口
  right=20mm,          % 切口
  headheight=15mm,     % 页眉高度
  footskip=15mm        % 页脚距离
}

% ==================== 字体设置 ====================
\setCJKmainfont[
  BoldFont=SimHei,
  ItalicFont=KaiTi
]{SimSun}

\setCJKsansfont{SimHei}
\setCJKmonofont{FangSong}

\setCJKfamilyfont{song}{SimSun}
\setCJKfamilyfont{hei}{SimHei}
\setCJKfamilyfont{kai}{KaiTi}
\setCJKfamilyfont{fang}{FangSong}

\newcommand{\song}{\CJKfamily{song}}
\newcommand{\hei}{\CJKfamily{hei}}
\newcommand{\kai}{\CJKfamily{kai}}
\newcommand{\fang}{\CJKfamily{fang}}

% ==================== 段落设置 ====================
\setlength{\parindent}{2em}        % 首行缩进2字符
\setlength{\parskip}{0pt}          % 段间距为0
\linespread{1.5}                   % 1.5倍行距

% ==================== 页眉页脚设置 ====================
\pagestyle{fancy}
\fancyhf{}
\fancyhead[CO]{\small\song 陈独秀最后的见解}
\fancyhead[CE]{\small\song 自由中国社丛书之二}
\fancyfoot[C]{\small\thepage}
\renewcommand{\headrulewidth}{0.4pt}
\renewcommand{\footrulewidth}{0pt}

% ==================== 标题格式设置 ====================
\ctexset{
  chapter = {
    format = \centering\huge\bfseries\song,
    name = {,},
    number = \arabic{chapter},
    beforeskip = 30pt,
    afterskip = 20pt,
  },
  section = {
    format = \centering\Large\bfseries\hei,
    name = {},
    number = \arabic{chapter}.\arabic{section},
    beforeskip = 20pt,
    afterskip = 10pt,
  },
  subsection = {
    format = \large\bfseries\kai,
    name = {},
    number = \arabic{chapter}.\arabic{section}.\arabic{subsection},
    beforeskip = 15pt,
    afterskip = 8pt,
  }
}

% ==================== 脚注设置 ====================
\setlength{\footnotemargin}{1em}
\renewcommand{\footnotelayout}{\small\song}

% ==================== 超链接设置 ====================
\hypersetup{
  colorlinks=true,
  linkcolor=black,
  citecolor=black,
  urlcolor=black,
  pdfborder=0 0 0,
  pdfcreator={XeLaTeX},
  pdfproducer={MiKTeX},
  pdftitle={陈独秀最后的见解},
  pdfauthor={陈独秀},
  pdfsubject={民主政治与自由},
  pdfkeywords={陈独秀, 民主, 自由, 政治}
}

% ==================== 自定义命令 ====================
\newcommand{\booktitle}{\Huge\bfseries\song 陈独秀最后的见解}
\newcommand{\bookauthor}{\Large\song 陈独秀\quad 著}
\newcommand{\publisher}{\small\song 自由中国社出版部发行}

% 引文环境
\newenvironment{quoteindent}
  {\list{}{\leftmargin=4em\rightmargin=0pt}\item\relax\fang\small}
  {\endlist}

% ==================== 文档开始 ====================
\begin{document}

% ==================== 封面+版权页（合并为一页） ====================
\thispagestyle{empty}
\vspace*{2cm}
\begin{center}
  {\booktitle}
  
  \vspace{1.5cm}
  {\bookauthor}
  
  \vspace{2cm}
  {\large\song 自由中国社丛书之二}
  
  \vspace{1.5cm}
  {\small\song 
    中华民国三十八年六月初版\\
    自由中国社出版部发行\\
    香港高士打道六十四号\\
    东南印务出版社印刷\\
    电话：二零八四八
  }
  
  \vfill
  {\small\song 版权所有\quad 翻印必究}
\end{center}

% ==================== 目录 ====================
\clearpage
\pagestyle{fancy}
\tableofcontents

% ==================== 正文 ====================
\mainmatter

% ---------- 序言 ----------
\chapter*{序言}
\addcontentsline{toc}{chapter}{序言\quad 胡适}
\markboth{序言}{序言}

\quad\quad 独秀是一九三七年八月出狱的，他死在一九四二年五月二十七日。最近我才得到藏的朋友们印行的《陈独秀的最后论文和书信》一小册，我觉得他的最后思想——特别是对于民主自由的见解，是"沉思熟了六七年"的，很值得我们大家想想。

独秀在一九三七年十一月给朋友的信里说：

\begin{quoteindent}
"我只注重我自己独立的思想，不迁就任何人的意见。我在此所发表的言论，已向人广泛的声明，只是我一个人的意见，不代表任何人。我已不属任何党派，不受任何人的命令指使，自作主张，自负责任。将来自食其果，现在完全不知道。我不怕孤立。"
\end{quoteindent}

\hfill ——给西流等的信

在那时候，人们往往还把他看作一个托洛斯基派的共产党人，但他自己在这里明白宣称他已"不属任何党派，不受任何人的命令指使"了。

一九三九年九月欧洲战事爆发之后，中国共产党在重庆出版的《新华日报》特别刊登反对一九一四年大战的论文，天天宣传此次战争是上次大战的重演，同是帝国主义者的战争。中国托派的《动向月刊》也这样说法。独秀很反对这种抄袭老话的论调，他坚决的主张：

\begin{quoteindent}
"赞助希特勒，或反对希特勒，事实上、理论上都不能含糊。反对希特勒，便不能同时打倒希特勒的敌人。否则所谓反对希特勒和阻止法西斯胜利，都是一句空话。"
\end{quoteindent}

\hfill ——一九四〇年三月二日给西流等的信

他更明白的说：

\begin{quoteindent}
"现在德俄两国的国家主义（纳粹主义及格别乌[G.P.U，按即秘密政治警察]政治，是现代的宗教法庭。此辈人类要前进，必须首先打倒比中世纪的宗教法庭还要黑暗的国家主义及格别乌政治。"
\end{quoteindent}

\hfill ——同年四月二十四日给西流等的信

那时候苏德互不侵犯条约已经订立，但斯大林对于希特勒的同情援助已很明显了。独秀在此时毫不迟疑的宣布他盼望世界大战的胜利属于英美。他说：

\begin{quoteindent}
"此次若是德国胜利了，人类将更加黑暗至少一个世纪；若胜利属于英法美，保持了资产阶级民主，然后才有道路走向大众的民主。"
\end{quoteindent}

\hfill ——同年给西流的信（在五六月之间）

他在这里提出了一个论断："保持资产阶级民主，然后才有道路走向大众的民主"——这个论断在一切共产党人的眼里是大逆不道的谬论。因为自从一九一七年俄国十月革命以来，共产党为了"无产阶级革命"的事业，造成了一套理论，说英美西欧的民主政治是"资产阶级的民主"，是资本主义的副产品，不是广大无产阶级所需要的民主。他们要打倒"资产阶级的民主"，要新建立"无产阶级的民主"。这是一切共产党人在那二十多年中说得烂熟的口号。托洛斯基失败之后虽然高喊党要民主，工会要民主，苏维埃要民主，但他并没有彻底想过整个政治民主自由的问题，所以"托派"的共产党人也都承袭了二十年来共产党攻击"民主"的烂调。在这一个重要问题上，列宁与托洛斯基、斯大林、希特勒与墨索里尼，是完全一致的，因为法西斯徒与纳粹徒都抄袭了国际共产党攻击"资产阶级民主"的老文章。

因此，独秀"保持资产阶级民主，然后才有道路走向大众的民主"的一句话，当时引起了他的朋友们一致的怀疑与抗议。一九四〇年七月独秀生病，只能勉强的答复。他说：

\begin{quoteindent}
"你们错误的根由，第一是不懂得资产阶级民主政治之价值（自列托以下均如此。）把民主政治看作是资产阶级的统治方式，是伪善，是欺骗，不懂得民主政治的内容是：

法院以外无捕人杀人权；

无政权不纳税；

非议会通过，政府无征税权；

政府之反对党有组织、言论、出版之自由；

工人有罢工权；

农民有耕种土地权；

思想、宗教自由；等等。

这都是大众所需要，也是三世以来大众以鲜血斗争七百余年，才得到今天的所谓'资产阶级的民主政治'。这是俄、意、德所要推翻的。

所谓'无产阶级的民主政治'，和资产阶级的民主只是实施范围的广狭不同，不是在内容上另有一套无产阶级的民主。

十月革命以来，拿'无产阶级的民主'这一个空洞的抽象名词做武器，来打毁资产阶级的实际民主，以至有今天的斯大林统治的苏联。意、德还是跟着学舌。现在你们又拿这一个空洞的名词做武器，来帮助希特勒攻打资产阶级民主的英美。"
\end{quoteindent}

\hfill ——一九四〇年七月三十一日给连根的信

（分段分行是我分的，目的是要醒目。）

这个简短的答复，是独秀自己独立思想的大觉悟。只有他能大胆的指摘"自列宁、托洛斯基以下均不会懂得资产阶级民主政治之价值"。只有他敢指出二十年来（现在三十年了）共产党用来打毁民主政治的武器——"无产阶级的民主"原来只是一个空洞的抽象名词！独秀的最大觉悟是承认"民主政治的内容"有一套最基本的信条——一套最根本的自由权利，都是大众所需要的，不是资产阶级所独有而大众所不需的。民主政治的内容，他在同一信的后文，做了一张对照表，如下：

\begin{center}
\small
\begin{tabularx}{\textwidth}{|X|X|}
\hline
\textbf{（甲）英美及战败前法国的民主制} & \textbf{（乙）俄德意的法西斯制} \\
& \small（原注：俄的政制是德意的老师，故可归为一类。）\\
\hline
（1）议会选举由各党（政府反对党也在内）自由竞争，选民有投票权。开会时有相当的讨论争辩。 & （1）苏维埃或议会选举均由政府指定。开会时只有举手，没有争辩。 \\
\hline
（2）法院外无捕人杀人权。 & （2）秘密政治警察可以任意捕人杀人。 \\
\hline
（3）政府的反对党至共产党公开存在。 & （3）一党专政，不容别党存在。 \\
\hline
（4）思想、言论、出版，相当自由。 & （4）思想、言论、出版，绝对不自由。 \\
\hline
（5）罢工本身非犯罪行为。 & （5）绝对不许罢工，罢工即是犯罪。 \\
\hline
\end{tabularx}
\end{center}

在这张对照表之后，独秀说：

\begin{quoteindent}
"每个共产主义者（按：独秀似不愿用'共产党'的名词，此处用'共产主义者'）看了这张表，还有脸说资产阶级民主不如无产阶级专政吗？宗教式的宣传时代，早已过去了，大家醒醒吧！今后的革命若仍主张'民主已过时，无产阶级政权只有独裁没有民主'，那只有使全世界的人类都堕入格别乌政治的黑暗地狱。"
\end{quoteindent}

独秀又说：

\begin{quoteindent}
"关于第二个问题（即民主政治制度问题），我研究俄国二十年来的经验，沉思熟了六七年，始决定了今天的意见。"
\end{quoteindent}

这是他自己的引言，下面的话，我现在摘引我最赞同的。在这里，他反复申论民主政治的重要，往往用十月革命以来的政制历史做例子。他说：

\begin{quoteindent}
"如果不实现大众民主，所谓'大众政权''无产阶级独裁'必然流为斯大林式的少数人的格别乌政治。这是事势所必然，非斯大林个人的心肠特别毒辣些。"
\end{quoteindent}

这是很忠厚的论断。向来"托派"共产党人把俄国的一切罪恶都归于斯大林一个人。独秀此时"已不属任何党派"了，所以他能看透托派的成见，指出俄国的独裁政制是一切黑暗罪恶的原因。独秀说：

\begin{quoteindent}
"斯大林的一切罪恶，是无产阶级独裁制之逻辑的发展。试问斯大林一切的罪恶，哪一样不是从一九一七年十月以来秘密的政治警察大权，党外无党，党内无派，不容思想、出版、罢工、暴动之自由，这一大串反民主的独裁制产生的呢？"
\end{quoteindent}

独秀自己加注道：

\begin{quoteindent}
"这些反民主的制度，都非创自斯大林。"
\end{quoteindent}

他又说：

\begin{quoteindent}
"若不恢复这些民主制度，斯大林之后的继承者，也不过是一个'专制魔王'。所以把俄国的一切坏事都归罪于斯大林，而不推源于制度之不良，仿佛只要去掉斯大林，一切都是好的，这种迷信个人轻视制度的偏见，公平的政治家是不应该有的。二十年的经验，尤其是近七八年的经验，使我们反省：我们若不从制度上寻出缺点，得到教训，只是闭起眼睛反对斯大林，将永没有觉悟。一个斯大林倒了，会有无数斯大林在俄国及别国产生。在十月革命后的俄国，明明是独裁制产生了斯大林，而不是有了斯大林才产生独裁制。"
\end{quoteindent}

独秀所主张恢复的民主制度，即是他所说的"民主政治之基本内容"。他在一九四〇年十一月写成《我的根本意见》一文，在本文内作一个更简明的总结：

\begin{quoteindent}
"民主主义是自从人类发生政治组织，以至政治消灭为止，各时代（希腊、罗马、近代以至将来）多数的人民反抗少数特权之旗帜。'无产阶级民主'，不是一个空洞名词，其具体内容也和资产阶级民主一样，要求一切公民都有集会、结社、言论、出版、罢工之自由。特别重要的是反对党派之自由。没有这些，议会与苏维埃同样一文不值。"
\end{quoteindent}

\hfill ——《我的根本意见》第八条

独秀在一年之内，前后四次列举民主政治的"内容"，这是最后一次，他看得更透彻了，所以能用一句话概括起来：民主政治只是一切公民（有产阶级的与无产阶级的，政府党与反对党）都有集会、结社、言论、出版、罢工之自由，更扼要一句：

\quad\quad 特别重要的是反对党派之自由。

在这三个字的短短一句话里，独秀抓住了近代民主政治制度的生死关头。近代民主政治与独裁政制的根本区别就在这里。承认反对党派之自由，才有近代民主政治。独裁制就是不容许反对党派的自由。

因为独秀"沉思熟了六七年"，认清了近代民主政治的本质内容，所以他能抛弃二十多年来共产党毁谤民主政治的成见，大胆的指出：

\begin{quoteindent}
"民主主义并非和资本主义及资产阶级是不可分的。"
\end{quoteindent}

\hfill ——《我的根本意见》第九条

他又指出：

\begin{quoteindent}
"近代民主制的内容比希腊时代丰富得多，实施的范围也大得多。因为近代是工业时代，我们便称之为资产阶级的民主制，其实这种制度不独为资产阶级所欢迎，而是广大民众流血斗争了五六百年才得到的。"
\end{quoteindent}

\hfill ——一九四〇年九月给西流的信

他很感慨的指出，俄国十月革命以后"轻率的把民主制和资产阶级统治一同推翻，以独裁代替了民主"，是历史上最可惋惜的一件大不幸。他说：

\begin{quoteindent}
"科学、近代民主制、社会主义，乃是近代人类社会三大天才的发明，至可宝贵。不料十月革命以来，轻率的把民主制和资产阶级统治一同推翻，以独裁代替民主，民主的基本内容被推翻了，所谓'无产阶级民主''大众民主'都是一些无实际内容的空洞名词。独裁制，无论以什么名义，其结果都只能是领袖独裁。任何独裁制都和残暴、蒙蔽、欺骗、贪污、腐化的官僚政治是不能分离的。"
\end{quoteindent}

\begin{quoteindent}
"布尔什维克人，更加把独裁制抬到天上，把民主骂得比狗屎不如。这种荒谬的观点，随着十月革命的胜利，征服了全世界。第一个采用这个观点的便是墨索里尼，第二个便是希特勒，首倡独裁制的本土——俄国，更是变本加厉，无恶不作。从此崇拜独裁的徒子徒孙普遍了全世界……"
\end{quoteindent}

\hfill ——同上

所以独秀"研究俄国二十年来的经验，沉思熟了六七年"的主要结论是：

\begin{quoteindent}
"民主之基本内容，无产阶级和资产阶级是一致的"；"民主主义并非和资本主义及资产阶级是不可分的"；"无产阶级民主"的具体内容也和资产阶级民主一样，要求一切公民都有集会、结社、言论、出版、罢工之自由，特别重要的是反对党派之自由；"任何独裁制都和残暴、蒙蔽、欺骗、贪污、腐化的官僚政治是不能分离的"。
\end{quoteindent}

\hfill ——《我的根本意见》第七、八、九条

以上是我摘抄的我的死友独秀最后对于民主政制的见解。在一九四一年一月九日他有给S和H的一封信，我引几句作为这篇介绍文字的结束：

\begin{quoteindent}
"弟自来立论，根据历史及现势之发展，而不迁就任何主义，更不喜引用前人之言论以为立论之前提。近作'我的根本意见'，亦未涉及任何主义。第七条主张重新估定布尔什维克的理论及其领袖（列宁、托洛斯基都包括在内）之价值，乃根据俄国二十年来的经验，非以马克思主义为尺度也。俄国革命的失败（成败不必计），即不合乎马克思主义，又安得而非之？予向不喜'圈子'，'圈子'即是'教派'。'正人君子'之于中国宋儒'道学'，此辈之于马克思主义，同一臭味也。弟偶见孔子之道有不对，便反对孔教；见得第三国际道理有不对，便反对他；对第四国际、第五国际……亦然。总之，兄是一个'终身的反对派'，是如此，然非弟故意如此，乃事实迫我不得不如此也。"
\end{quoteindent}

因为他是一个"终身的反对派"，所以他不能不反对独裁政治，所以他从苦痛的经验中悟得近代民主政治的基本内容，"特别重要的是反对党派之自由。"

\hfill 一九四九年四月四夜，在太平洋船上

% ---------- 书信部分 ----------
\chapter{给西流等的信（一九四〇年三月二日）}
\markboth{给西流等的信}{一九四〇年三月二日}

\quad\quad 和韵兄等：

以上的话，是因为十月革命得以保全。现在呢，我们也不该抄袭列宁一九一四年大战时的理论。一切理论与口号都有其时代性与特殊性，是不能随便抄袭的。对于像欧洲大战这样重大的事件，不能忽视其时代的特点，认为是历史重演，只背诵一大篇空洞的理论了事，真正的马克思理论家是不抄袭老话的啊！历史不会重演，错误是会重演的。上次大战时，一切帝国主义都是为了奴役本国人民和掠夺殖民地而战。《动向月刊》竟做了他们的应声虫，在这一点上，我实在看不出中国托派与斯大林派之区别。列宁当年对一九一四年大战的理论之正确，是由于不抄袭恩格斯、马克思之现成的理论，而是自己用脑子观察分析当时欧洲大战的环境特点；其口号之收效，是由于俄国是被压迫民族，而且德国不能加以布列斯特和约的压迫。在重庆的《新华日报》，天天大登其列宁反对一九一四年大战的文章，天天说此次战争是上次大战的重演，希特勒政权垮台了，他们却一面在国际上站到希特勒方面，一面宣传帝国主义大战，促成英国工人反对战争，英国共产党四人，因赞成希特勒被开除，第三国际也是帮助他取得胜利。

没有"组织国际联合反法西斯"的口号，他们却在国内用不通的"人民阵线""反侵略"口号，做勾结希特勒政府的梦，而不是组织国际普遍反法西斯的联合。等到法国贝当政府和法西斯德国合作之后，他们才改变论调，这不是太晚了吗？兄来信也说："如果是法西斯胜利，人类将堕入黑暗，因此必须阻止法西斯的胜利。"这话对了，但阻止法西斯胜利的方法是什么呢？我只看见你们一方面反对英国政府，一方面又不主张联合英国反抗法西斯，这不是自相矛盾吗？

反对希特勒，便不能同时打倒希特勒的敌人，否则所谓反对希特勒和阻止法西斯胜利都是一句空话。如何才能阻止法西斯胜利呢？只有联合希特勒的敌人，共同反抗希特勒。今天的武器和交通都和以前大大地不同了，英国的国内革命即使能够成功，在希特勒的强大军事力量面前，也难以持久。只有希特勒对英国失利，像以前拿破仑第三一样，引起国内革命，才能阻止法西斯胜利。若在英国国内发动革命，推翻现政府，只能使英国陷入混乱，反而帮助希特勒征服欧洲，这不是反动是什么？

历史不会重演，但错误是会重演的。有人把一九一四年大战的理论口号用于今日战争，而忘记了被压迫民族的特点，无论他唱得如何高调，只能有助于日本；现在又有人把列宁当年的理论用于此次欧洲战事，而忽略了反法西斯的特点，无论他说得如何动听，只能有助于希特勒。英国虽然不是被压迫的普鲁士，但希特勒却是要征服全欧洲的新的第三帝国，而不是威廉第二；因此，不但在德国，就是在全欧洲，反法西斯的斗争都是被压迫民族的解放斗争。

今天我们的任务，是联合一切反法西斯的力量，共同打倒希特勒及其同伙，而不是重复过去的错误，把民主国家和法西斯国家混为一谈。如果我们仍然坚持"打倒一切帝国主义"的空洞口号，而不区别法西斯与一般帝国主义的不同，只能是帮助法西斯，危害人类的解放事业。

\hfill 一九四〇年三月二日

\chapter{给西流等的信（一九四〇年四月二十四日）}
\markboth{给西流等的信}{一九四〇年四月二十四日}

\quad\quad 西流兄等：

前函有未尽之意，再补充如下。我有两个信念：（一）在此世界大战之前，甚至之后短期内，中国的民主革命无胜利之可能。（二）在德意日法西斯格别乌政治大胜利的时代，一切（反帝也包括在内）都比打倒法西斯格别乌政治次要。因此，一切不利于这一斗争的，都是反动的。我的以上见解，不但在英法美国内反对法西斯是反动的，即印度独立运动如果不利于反法西斯斗争，也是反动的。民族解放运动固然是进步的，但如果在大战期间，印度一旦独立，必然陷入日本的统治，使希特勒对英国取得决定性的胜利，这不是反动是什么？我这一意见，不但听起来奇怪，兄等亦必认为荒谬，因为这和我们中国以前所信奉的公式太冲突了。兄等如能反省一下，抛弃以前抄袭来的公式，那便更好。

反对国家主义及格别乌政治的斗争，不是由于民族成见，而是由于这种政治是人类进步的最大障碍。即使英美等民主国家有种种缺点，我们也应该支持他们，因为他们的民主制度是反法西斯的基础。若高唱"打倒一切帝国主义"，使民主国家失败，那更是罪恶滔天。

\hfill 一九四〇年四月二十四日

\chapter{给连根的信（一九四〇年七月三十一日）}
\markboth{给连根的信}{一九四〇年七月三十一日}

\quad\quad 连根兄：

你们的意见我都见了，不得不力疾给你详细答复：你们错误的根由，第一是不懂得资产阶级民主政治之价值（自列宁、托洛斯基以下均如此。）把民主政治只是看作资产阶级的统治形式，是伪善，是欺骗，而不懂得民主政治的内容是：法院以外无捕人杀人权；无政权不纳税；非议会通过，政府无征税权；政府之反对党有组织、言论、出版自由；工人有罢工权；农民有耕种土地权；思想、宗教自由；等等。这都是大众所需要，也是三世以来大众以鲜血斗争七百余年才得到今天的所谓资产阶级民主政治，这是俄、意、德所要推翻的；所谓"无产阶级的民主政治"，和资产阶级的民主只是实施范围的广狭不同，不是在内容上另有一套无产阶级的民主。十月革命以来，拿"无产阶级的民主"这一空洞的抽象名词做武器，来打毁资产阶级的实际民主，以至有今天的斯大林统治的苏联，意、德还是跟着学舌。现在你们又拿这一个空洞的名词做武器，帮助希特勒攻打资产阶级民主的英美。

\hfill 一九四〇年七月三十一日

\chapter{给西流的信（一九四〇年九月）}
\markboth{给西流的信}{一九四〇年九月}

\quad\quad 西流兄左右：

日前寄上一信，内论欧洲战事及民主问题，想已收到。七月三十一日手示及守一兄的信已悉。因病不能早复，今犹如此。（此次害了日疟，精神不佳，请勿多疑！）

来函说明："你对民主的了解，和对于世界的局势观察，我觉得还不免有些'头巾气'"，我所争论的中心，正是这两点：（一）大国是否有革命局势。（二）是否应该保卫民主。你既然说我有"头巾气"（其实是反动），又说我没有错，即你自己也感觉到有自相矛盾之处吧！

关于第一个问题，我只能答复一个"否"字。无论在法国，还是在英国，资产阶级和希特勒相比，我无疑的否定英、法会有革命的局势，其理由是：（一）两国的革命力量，已被斯大林派和改良派消灭殆尽；（二）各国的资产阶级有了一八七一年巴黎公社和一九一七年俄国革命的教训，采取了许多让步政策，把武器完全掌握在自己手中，防止国内革命；（三）此时德国的武器和交通及统治征服地的方法，均非一八七一和一九一七可比，英法政府军队战败后，民众一时决不能突起反抗；（四）德国未得世界霸权之前，一旦采取守势即可巩固其统治，失败后也不会仍保持法西斯政权，此种情况与法国大革命时期不同，民众及其他自由派会抬头，然而这只能有利于民主革命之开始，很难立即走上社会主义革命，更不能以无产阶级专政做号召，实际上因为以前我们所倡导的"帝国主义大战后失败国将引起革命"这一公式完全被推翻了，只有迷信公式的人才会仍然坚持这种观点。

兄能受得了这种结论吗？如此，必须讨论第二个问题，即我们中间主要的不同意见还是在于"是否应该保卫民主"。一个格别乌式的俄国已经足够使人窒息了，再加上格别乌式的德国、意大利，你老兄不觉得人类太可怜了吗？有人说"民主与独裁没有区别，不管谁胜谁负，只得听天由命"，这话只能取决于战争之结果：如果法西斯胜利，即使英国失败引起了革命，也只有使世界更加黑暗；如果民主国家胜利，至少可以保持一部分民主自由，为将来的进步奠定基础。

民主与独裁像一所房子：不管好坏，只得取决于法西斯与民主独裁之斗争结果。英美是民主国家，德国是独裁国家，帮助希特勒征服了有民主传统的法国，我们可以进一步说，如果人们仍然轻率的认为"白里安和希特勒是一样的"，因此帮助希特勒取得政权；现在纳粹德国和民主的英法是一样的，又帮助希特勒征服欧洲，人类还有希望吗？在英国采取失败主义，除了帮助希特勒胜利之外，还有什么别的结果？历史事实证明，你们的错误在于混淆了民主与独裁的界限，把反对资产阶级政府和反对法西斯独裁混为一谈。

我认为：以大众民主代替资产阶级的民主是进步的；以德俄的独裁代替英美的民主，是退步的。直接或间接有意无意的帮助这一退步的人，都是反动的，不管他口中说得如何左。我还认为：所谓"大众政权"如果不能实现大众民主，便必然流为斯大林式的少数人的格别乌政治，这是事势所必然，非斯大林个人的心肠特别毒辣。

关于第二个问题，我研究俄国二十年来的经验，沉思熟了六七年，始决定了今天的意见。（一）民主不是一个抽象名词，有它的具体内容，资产阶级的民主和无产阶级的民主，其内容大致相同，只是实施的范围不同而已（见前信及对照表）。（二）民主之内容固然包含议会制度，而议会制度不是民主的全部内容。许多年来，许多人把民主和议会制度看作一件事，排斥议会制度，同时便排斥民主，这是大错特错。议会制度可以成为历史陈迹，民主则不然，如果排斥议会制度而没有民主内容，仍是一种形式主义的民主代议制，甚至俄国的苏维埃，比资产阶级的形式民主议会还不如。（三）民主是自古代希腊以至今天乃至将来，每个时代被迫的大众反抗少数特权的旗帜，并非只是某一特殊时代的历史现象，并非是资产阶级统治的形式。如果民主已经过时，一去不复返，同时便可以说政治及国家也已过时即已死亡了。如果民主只是资产阶级的统治形式，无产阶级的政治形式只有独裁没有民主，则斯大林所做一切罪恶都是应该的了，列宁所说"民主是对于官僚制的抗毒素"，乃成了一句空话。主张工人及民众的民主，也是等于叫天不应，等于老百姓向暴君乞求恩赐，而不是通过斗争取得。如果无产阶级民主和资产阶级民主不同，那是完全不了解民主之基本内容（法院外无捕人杀人权，政府反对党公开存在，思想、言论、出版、罢工之自由权利），无产阶级和资产阶级是一样需要这些民主权利的。

如果斯大林的罪恶与无产阶级独裁制无关，即是说斯大林的罪恶非由于十月革命以来所建立的制度违反了民主制之基本内容（这些反民主的制度，都非创自斯大林），而是由于斯大林的个人心肠特别坏，这完全是唯心派的见解。斯大林的一切罪恶，乃是无产阶级独裁制之逻辑的发展。试问斯大林一切罪恶，哪一样不是从十月革命以来秘密的政治警察大权，党外无党，党内无派，不容思想、出版、罢工、暴动之自由，这一大串反民主的独裁制产生的呢？若不恢复这些民主制度，斯大林之后的继承者，谁也不过是一个"专制魔王"。所以把俄国的一切坏事都归罪于斯大林，而不推源于制度之不良，仿佛只要去掉斯大林，一切都是好的，这种迷信个人轻视制度的偏见，公平的政治家是不应该有的。二十年的经验，尤其是近七八年的痛苦经验，使我们反省：我们若不从制度上找出缺点得到教训，只是闭起眼睛反对斯大林，将永没有觉悟。一个斯大林倒了，会有无数斯大林在俄国及别国产生。在十月革命后的俄国，明明是独裁制产生了斯大林，而不是有了斯大林才产生独裁制。如果认为民主已经过时，不必恢复民主，即等于说无产阶级政权不需民主，这将使人类永远堕入黑暗！

（四）近代民主制的内容比希腊时代丰富得多，实施的范围也大得多，因为近代是工业时代，我们便称之为资产阶级的民主制，其实这种制度不独为资产阶级所欢迎，而是广大民众流血斗争了五六百年才得到的。科学、近代民主制、社会主义，乃是近代人类社会三大天才的发明，至可宝贵。不料十月革命以来，轻率的把民主制和资产阶级统治一同推翻，以独裁代替民主，民主的基本内容被推翻了，所谓"无产阶级民主""大众民主"都是一些无实际内容的空洞名词。独裁制像一把利刃，今天用之别人，明天便会用之自己。列宁当时也曾经警觉到"民主是对于官僚制的抗毒素"，而亦未采用民主制，如取消秘密政治警察，容许反对党派公开存在，思想、罢工自由，直至这把利刃伤害到他自己，想到要党内民主、工会和各苏维埃民主自由，然而太晚了！其余一班无知的布尔什维克人，更加把独裁制抬到天上，把民主骂得比狗屎不如。这种荒谬的观点，随着十月革命的胜利，征服了全世界。第一个采用这个观点的便是墨索里尼，第二个是斯大林，第三是希特勒，这三个反动堡垒，把现代世界变成了新的中世纪，他们把有思想的人变成无思想的机器牛，任其鞭打驱使，人类若无力推翻这三个反动堡垒，只有做机器牛的命运。所以目前全世界的一切斗争，必须以推翻这三大反动堡垒为首要任务，否则任何好听的名词，如无产阶级革命、民族革命，都无非是在客观上帮助这三大反动势力。如果我们承认当前推翻这三大反动堡垒是首要斗争任务，第一必须承认即英法美不彻底的民主制也有保卫的价值；第二必须抛弃刘仁静、彭述之的理论，即：任何时期，任何事件，无产阶级都不能与资产阶级共同行动；这一理论，固然在北伐战争中，在抗日战争中，都不能应用，在目前国际斗争中也同样不能应用，若采用这一理论，都只有反动作用。

附上对照表一张，以供兄参考：

\begin{center}
\small
\begin{tabularx}{\textwidth}{|X|X|}
\hline
\textbf{（甲）英美及战败前法国的民主制} & \textbf{（乙）俄德意的法西斯制} \\
& \small（原注：俄的政制是德意的老师，故可归为一类。）\\
\hline
（1）议会选举由各党（政府反对党也在内）自由竞争，选民有投票权。开会时有相当的讨论争辩。 & （1）苏维埃或议会选举均由政府指定。开会时只有举手，没有争辩。 \\
\hline
（2）法院外无捕人杀人权。 & （2）秘密政治警察可以任意捕人杀人。 \\
\hline
（3）政府的反对党至共产党公开存在。 & （3）一党专政，不容别党存在。 \\
\hline
（4）思想、言论、出版，相当自由。 & （4）思想、言论、出版，绝对不自由。 \\
\hline
（5）罢工本身非犯罪行为。 & （5）绝对不许罢工，罢工即是犯罪。 \\
\hline
\end{tabularx}
\end{center}

守一对大战的见解，是由于估计苏联性质及民主价值出错而产生的，我与之相反，但我们都是出于真诚的，惟有你老兄对大战的意见同意于守一等，对苏联与民主似乎还是和我接近，这实在不可理解。此信请抄送守一及其他朋友。原信及各附件，我暂保留，因为打算将来印行，供后人参考。昌兄的信附上，请你转交。顺祝健康！

\hfill 弟仲白\quad 一九四〇年九月

\chapter{给四流的信（一九四〇年九月）}
\markboth{给四流的信}{一九四〇年九月}

\quad\quad 昌兄：

你说"在战争进行中之现在，民主与法西斯之界限已然消失，或将消失"。这句话真是莫名其妙！（一）在政制本身上，民主与法西斯有本质的不同，这种界限永远不会消失；（二）其所谓界限消失，是指英法美等民主国家日渐法西斯化；即真是如此，也绝对不能据此以作为我们欢迎独裁反对民主的理由；（三）英法美将来是否法西斯化，是要靠第三国际、第四国际帮助希特勒完全胜利；希特勒军队打到英国本土，才能实现，否则英法美的民主传统不是轻易可以推翻的；（四）为现在的民主国家采取一些加强统治的措施是反法西斯的需要，便不懂得法西斯究竟是什么。

你把民主国家的战时措施看作是法西斯化，这是错误的。民主国家的战时措施是为了保卫民主，而法西斯的统治是为了建立独裁，二者有本质的区别。如果把英国的内阁改组、加强军事力量看作是法西斯化，那便是混淆了民主与独裁的界限。

守一对大战的见解，是由于估计苏联性质及民主价值出错而产生的，我与之相反，但我们都是出于真诚的，惟有你老兄对大战的意见同意于守一等，对苏联与民主似乎还是和我接近，这实在不可理解。希望你能睁大眼睛看看左列对照表，便知道民主与法西斯的界限并没有消失，也不会消失。

\begin{center}
\small
\begin{tabularx}{\textwidth}{|X|X|}
\hline
\textbf{（甲）英美及战败前法国的民主制} & \textbf{（乙）俄德意的法西斯制} \\
& \small（原注：俄的政制是德意的老师，故可归为一类。）\\
\hline
（1）议会选举由各党（政府反对党也在内）自由竞争，选民有投票权。开会时有相当的讨论争辩。 & （1）苏维埃或议会选举均由政府指定。开会时只有举手，没有争辩。 \\
\hline
（2）法院外无捕人杀人权。 & （2）秘密政治警察可以任意捕人杀人。 \\
\hline
（3）政府的反对党至共产党公开存在。 & （3）一党专政，不容别党存在。 \\
\hline
（4）思想、言论、出版，相当自由。 & （4）思想、言论、出版，绝对不自由。 \\
\hline
（5）罢工本身非犯罪行为。 & （5）绝对不许罢工，罢工即是犯罪。 \\
\hline
\end{tabularx}
\end{center}

每个共产主义者看了这张表，还有脸说民主与法西斯没有界限吗？还有脸说资产阶级民主不如无产阶级独裁吗？宗教式的宣传时代，早已过去了，大家醒醒吧！今后的革命若仍主张"民主已过时，无产阶级政权只有独裁没有民主"，那只有听任格别乌政治统治全人类，还有什么民主可言？

\hfill 一九四〇年九月

\chapter{我的根本意见（一九四〇年十一月二十八日）}
\markboth{我的根本意见}{一九四〇年十一月二十八日}

\quad\quad （一）必须抛弃"人民革命"的空想。力量越大阻力也越大这一物理现象，同样也可以应用于革命。统治阶级掌握着国家机器，拥有强大的军事力量和经济力量，要推翻他们的统治，必须有足够的力量做基础，否则革命必然失败。我们不能把革命看作是轻而易举的事情，更不能把反动的势力看作是革命的局势。

（二）无产阶级的革命，不会在任何时间、任何国家，都有革命局势。最荒谬的是把反动的局势，看作是革命局势。尤其是在大众遭到重大失败之后，社会经济陷入大恐慌之时，无产阶级的力量受到严重打击，很难发动革命。

（三）此次世界大战，自然不是单纯的资本帝国主义之间的战争，而是发展到第二阶段之开始，即是由多个帝国主义的国家，合并为简单的两个对立的帝国主义集团之间的战争。不能把无产阶级的力量估计过高，我们不能轻率宣布"资本主义已到末日"，没有发动全世界无产阶级革命的条件。

（四）严格区分小资产阶级"中间派"的武断性，和无产阶级"集中制"的自然性之间的不同。小资产阶级中间派往往脱离实际，提出一些不切实际的口号，而无产阶级的集中制是为了团结力量，实现革命目标。

（五）严格区分小资产阶级分子和坚定的无产阶级分子之间的不同。小资产阶级分子往往动摇不定，容易受反动势力的影响，而坚定的无产阶级分子能够坚持革命原则，为实现革命目标而奋斗。

（六）现在不是最终革命时代，不但在落后国家，即在欧美先进国家，如果有人武断的说：小资产阶级已经没有一点进步作用，已完全走到反动的道路，这只是下了一个错误的判断，将来必然会遭到事实的反驳。

（七）应该无成见的觉悟俄国二十年来的教训，科学的而非宗教的重新估定布尔什维克的理论及其领袖（列宁、托洛斯基都包括在内）之价值，不能把一切错误都归于斯大林一人，例如无产阶级政权之下民主制的问题。

（八）民主主义是自从人类发生政治组织，以至政治消灭为止，各时代（希腊、罗马、近代以至将来）多数的人民反抗少数特权之旗帜。"无产阶级民主"不是一个空洞名词，其具体内容也和资产阶级民主一样，要求一切公民都有集会、结社、言论、出版、罢工之自由。特别重要的是反对党派之自由。没有这些，议会与苏维埃同样一文不值。

（九）政治上的民主主义和认识上的实证主义，是相成而非相反的东西。民主主义并非和资本主义及资产阶级是不可分的。无产阶级政权因反对资产阶级及资本主义，遂并民主主义而亦反之，即各国所谓"无产阶级革命"出现了，而没有民主制做官僚制之抗毒素，也只是世界上出现了一些斯大林式的官僚政治，贪污、腐化、堕落，不能实现社会主义的理想。

（十）"无产阶级独裁"，根本没有这种东西，即党的独裁，结果也只能是领袖独裁。任何独裁制都和残暴、蒙蔽、欺骗、贪污、腐化的官僚政治是不能分离的。

（十一）此次国际大战，自然是两个帝国主义集团争夺全世界霸权的战争。所谓"为民主自由而战"自然是一种宣传口号；然而不能因此便否认英美民主国家有若干民主自由之存在。在那里，在野党、工会、罢工之存在，是客观事实，除了纳粹、法西斯的爪牙，是不能用任何理由来否认的。我们更未听到美国用纳粹对待犹太人的方法来对待孤立派。希特勒的徒子徒孙，以其统治德国的野蛮黑暗的方法统治全世界，是以比中世纪宗教法庭还要野蛮黑暗的法治统治全世界，使全世界只有一个主义，一个政党，一个领袖，不容任何异己之存在，并不容被征服的国家中土著纳粹及各色各样的土法西斯之存在。希特勒之胜利，将使全人类由有思想有精神自由的人，变成无思想无自由意志的机器；所以世界各国（当然包括中国在内）有良心的进步分子，在此大战一开始以及现在将来，都应以"消灭希特勒的纳粹徒"为各民族共同攻击之目标，其他一切，只有属于这一总目标有正面的作用而非负面的作用。有进步思想的人，若因希特勒的纳粹一胜利，社会主义，甚至民主主义，民族解放，一切都无从谈起。

（十二）在此帝国主义大战中，对民主国家方面采取失败主义，采取以国内的革命代替反对帝国主义战争的策略，无论口号说得如何好听，事实上只有帮助法西斯，例如英国人自己采取失败主义，必然会使英国陷入混乱，帮助希特勒征服欧洲。

（十三）战争与革命只有在趋向进步的国家，是生产力发展的结果，又转而成为生产力发展的原因；若在反动的国家，反而使生产力更加削弱，使民众品格更加堕落，奢侈、腐化、自私自利。

（十四）在纳粹统治之下，更没有任何理由可以说：革命之爆发，在有民主成份的政权之下比在独裁政权之下更容易；也没有任何理由可以说：纳粹胜利比其失败对德国革命更有利。我们没有理由相信在人类自由之命运上斯大林党徒比希特勒党徒好多少。

（十五）帝国主义以殖民地、半殖民地存在为条件，正如资本主义以私有财产为存在条件。我们不能幻想资本统治不崩溃可以取消私有财产制，同时也不能幻想殖民地、半殖民地的民族独立运动，不和帝国主义国家（宗主国及宗主国的敌对国家）中的革命运动结合起来，会得到胜利。在今天——英美和德国、日本两大帝国主义争夺世界霸权的今天，孤立的民族运动，无论由何阶级领导，不是完全失败，便是更换主人，或者还是更换一个更凶恶的主人；即使更换一个较开明的主人，也不利于民族的真正解放和发展。

\hfill 一九四〇年十一月二十八日

\chapter{给S和H的信（一九四一年一月十九日）}
\markboth{给S和H的信}{一九四一年一月十九日}

\quad\quad S、H二先生：

与先生别三年矣，与S先生更二十余年不见了。回忆北京往事，不胜感慨！

见二先生的来信，对我的近作有所指教，感激之至。我向来立论，根据历史及现势之发展，而不迁就任何主义，更不引用前人之言论以为立论之前提，这种"不迷信权威"的态度，是科学的态度，而非宗教的态度。

总之，逃避现实与迁就现实，均经不起事实的考验及时代的淘汰，我不取之。言论的长短，不问其出自何家，若内容荒谬，因其出品而崇拜之，是错误的；若内容正确，因其出自反对派而排斥之，也是错误的。吃肉，只要味道好，不管它来自何家店铺；探求真理，只要道理正确，不管它来自哪个学派。

我对于中国宋儒"道学"，向来不感兴趣。我见得孔道有不对之处，便反对孔教；见得第三国际道理有不对之处，便反对第三国际；见得第四国际、第五国际……道理有不对之处，亦然。总之，我是一个"终身的反对派"，并非我故意如此，乃事实迫我不得不如此也。

近作《我的根本意见》，亦未涉及任何主义，第七条主张重新估定布尔什维克的理论及其领袖（列宁、托洛斯基都包括在内）之价值，乃根据俄国二十年来的经验，非以马克思主义为尺度也。若布尔什维克的理论及领袖的主张正确，（成败不必计）即不合乎马克思主义，我也不能非难之。"圈子"即是"教派"，"正统"等说法，都是宗教式的偏见，非科学的态度。

今后如有新作，自当寄上乞教，无奈多病不能多作，即油印亦不易办到也。此祝健康！

\hfill 弟独秀\quad 一九四一年一月十九日

\chapter{战后世界大势之轮廓（一九四二年二月十日）}
\markboth{战后世界大势之轮廓}{一九四二年二月十日}

\quad\quad 历史决不会重演，此次大战已使世界各方面发生巨大变化，或已产生巨大变化之萌芽，拿过去的理论公式来硬套将来的历史事实，简直不中用了。

此次大战不外三种结果：一是英美和德日不分胜负而讲和；二是胜利属于英美；三是胜利属于德日。第一种结果之可能最少，我们似不必加以讨论。第二种和第三种结果以何者最大呢？以现势观之，自然以德日占优势。开战已两年多了，又因得到苏联的支持，英国全然休息了一年，此次以全力反攻仍不能击退德国在北非的少量军队，若能够于最近的将来击败德国的大军，这是很难以想象的事。若英国在各战场之失败，都由于军事力量之不足，再过一年或至多一年半，美国武器生产大规模扩大之后，战局便可能改变；有人高呼"全面反攻"，但过去直至现在，政府官吏之因循守旧以及工厂主只顾私人利益，导致武器生产之缓慢，英美能否胜过德国本土及其可能利用之附庸国，确是一个大问题。即假定将来英美可以胜利，又有何种力量能够使德国及其伙伴，在此一年至一年半之内，按兵不动，等待美军准备就绪呢？德国的内部危机，固然为英美所期望，然而在对外战争取得胜利的情况下，未必会爆发。德国唯一的弱点是缺乏煤油，也只是在她始终无力夺取高加索或伊朗的条件之下不能够坚持长期战争。总之，双方的情况各有优劣，胜利的关键在于能否争取时间，扩大军事力量。

第一：所以迫在眉睫的德国的春季攻势，无论在北非、中东或直捣莫斯科，其胜负都可说是决定此次大战全局成败之最大关键。一旦德日取得胜利，英美是不能长久支持下去的。自古至今，战争是力量的较量，不是单纯的正义与否可以决定胜利的重要因素。

若德日胜利，英美失败，世界便会形成德日两大霸权国家对立的局面。英国失去欧洲、北非、近东以至中东，已非易事，一时决无力顾及远东以至南洋、澳洲，这些地区自然属于日本的势力范围。德国将统治欧洲大陆，决定欧洲的命运。

若胜利属于希特勒，英国便完了，美国也只能退守美洲以自保。希特勒胜利之后，德国仍然要向西方扩张，自印度以东，则非他的军事力量直接所能及。那时无论日本是否与德国成立正式同盟，日本都将是美国的主要敌人。美国固然未必立即对日本宣战，希特勒在未征服美国以前，他也不会为了德国的利益而开罪日本，强迫他的盟友向美国宣战，自大西洋攻击美国。希特勒知道在军事力量上，日本得到的利益而美国受到的损失，对德国有利。

世界还会有大战，我们还不能知道，但所能知道的只是在战争的因素未除去以前，战争的结果是不可避免的。而且若德国胜利，下次大战必然是德美之间的战争，其规模将比此次大战更大，破坏性更强。

美德之间虽然无直接冲突，但国际利益的争夺必然会导致战争。希特勒征服欧洲之后，必然会向美洲扩张，与美国争夺世界霸权。美国为了保卫自己的利益，也必然会奋起反抗，这是事势所必然。

改良资本主义制度既非易事，消灭资本主义制度更不能够如人们所想像的那样轻易。此次大战，不但英美，就是在德日国内，也必然企图改良资本主义制度，以巩固其统治。希特勒一向攻击资本主义，但他并不能消灭资本主义，只是他和他自己开开玩笑。他们改良的方法，不外：拿协定以至同盟，调整各集团内的关系；用物物交换减少各集团内的贸易摩擦；拿部分国有化，来代替某些私人垄断。集团内的关系即使能够调整，对集团外的关系必然要更加紧张；物物交易不但不能全部实行，即所能够实行的，仍以货币计算，仍是商品交换，而非分工互助。某国部分国有化，已是前世就有的事，全部国有化，即实行所谓国家资本主义，理论上好像说得通，而事实上必不可能。掌握生产工具的大资本家，不会自愿的把私有财产交给国家，这是不可想像的事。如果有人幻想能有一个超然政府来和平的没收资本家的财产，那更是空想。所谓"超然政府"，上午成立，下午即被资本家收买了。所以以上三种改良方法，决不能动摇资本主义制度的基础。资本主义制度一旦开始产生，便必然会随着自身的发展而日益强大，一切改良方法都不能阻止其内在矛盾的发展，反而只有使矛盾更加尖锐，最终导致危机爆发。

为了必须把国内生产过剩的产品推向国外市场，必须向外掠夺原料，便不得不加高关税壁垒，实行贸易保护主义，以至发动战争。这一连串因果相连的现象都是资本主义制度的必然产物，决不是思想、道德所能改变的。在世界强大的帝国主义国家争夺市场、原料的今天，任何国家要想独善其身，都是不可能的。

资本主义制度的基本规律是：生产力发展，追求利润最大化，导致生产力与购买力之间的矛盾，因生产过剩，而物价低落，而工人失业，而工资下降，而形成恐慌；经过一个时期，因生产力停滞而恢复平衡，因生产力发展而再次出现危机，如此恐慌之周期循环不已。自发生生产过剩之危机以来，解决的方法有二：一是自助的销毁少量以至大量生产过剩产品，这是愚蠢可笑的办法；二是疯狂的掠夺殖民地及国外市场，这是反动的办法。这两种办法都不能从根本上解决问题，只能使危机更加严重，最终导致战争。

第一次世界大战之后，威尔逊的"十四点计划"，本想建立国际联盟，维护世界和平，但结果却失败了，反而导致了第二次世界大战的爆发。这说明在资本主义制度之下，没有真正的和平，只有暂时的休战。此次大战，又有人幻想战后建立"世界政府"，维护世界和平，这同样是不切实际的空想。在资本主义制度之下，各国为了争夺利益，必然会不断发生冲突，战争是不可避免的。

日本和苏联，当然都有各自扩张其势力范围之野心。然而生产力的发展最终要决定他们的命运。其他殖民地及半殖民地国家，若想脱离帝国主义的统治，建立新的独立国家，其时代已经过去了。在各集团范围内，国家力量之强弱决定其地位。我们略分四类：第一等是较有面子的所谓"同盟国"，例如日本之于德国，苏联之于英国；第二等是附属国，例如意大利之于德国，荷兰、比利时之于美国，虽然有自己的政府，但内政外交都受大国控制；第三等是被保护国，例如法国、比利时之于德国，丹麦、意大利之于英国，菲律宾之于美国，虽然有一个自己的政府，而不能有独立之外交；第四等是殖民地，连自治政府也没有，统治权操在宗主国总督之手。比殖民地更低一等的，自来是没有的，有之便是孤立民族日渐消灭之美洲、澳洲的土人。各集团圈内的国家民族地位高低不同，而有一共同之点，即是他们的政治及社会制度都必须或多或少的按照领导国的模样改造，根本相反的制度是不能够存在的：德国所领导的集团，多少都要实行纳粹制度；美国所领导的集团内，多少也要按照民主制度进行；苏联所领导的集团，自然也要实行苏维埃制度。这要看领导国的命运及成功与否，能够影响整个集团。俄国革命的经验证明，帝国主义世界中最薄弱的一环之突破，不能够使整个帝国主义制度瓦解。至于现在的苏联，不但她的生产力不能胜任领导世界革命，自身早已离开了社会主义的道路，走上了国家资本主义的道路。

有些喜做梦的人，当此大战一开始，便梦想弱小民族独立的机会到了：其实欧洲的殖民地，一旦英美失败便落入了日本的掌握；非洲的殖民地一旦英国失败便落入了德国的掌握。有人甚至想战争会引起的社会主义革命就快到来，不幸幻想破灭了，他们的悲哀，已不可言状；如果再觉悟到今后民族运动都会受到限制，而且纳粹若统治整个地球，人们将感觉到由喜悦的天国坠入到悲哀的深渊，感觉到命运注定要走下坡路了。其实人类文化史，始终是沿着曲折的道路前进的，此次战争既非意图建立天国，也不是要使人类坠入毁灭的深渊，人们对自由幸福的希望，不应因一时的成败而感到失望或悲哀，历史的进步是不可阻挡的。

在黑暗中求进步的例子，只有眼光狭隘的党派之人才看不到这一点。历史也和地球一样，无论在光明的白昼或黑暗的深夜都是运行不息的。此次大战，即使不幸得很，胜利果然属于德日，竟至统治了地球，占人类大多数的人民，在政治上将受到长期窒息的大灾难，而在经济上，德日的胜利，固然不能动摇资本主义制度的生存，反而会加强物质文明的发展，这本来是资本主义在血腥的战争中发展的规律。例如德国的军事工业，比战前有了很大的发展；日本的工业也有了一定的进步。而且"民族运动会受到限制"这句话，不一定完全是坏事。无论是全世界统一为一个国家，没有革命的统一，反革命的统一也有进步的意义。例如中国历史上的秦始皇统一中国，虽然是封建专制的统一，但也促进了经济文化的发展。从全世界的进步来看，民族界限的消除，是历史发展的必然趋势。

真正的民族解放，只能和帝国主义国家的社会主义革命同时实现；在资本主义帝国主义时代，任何弱小民族，若想依靠自身的力量，排除一切帝国主义之干涉，以实现孤立的民族独立，是没有出路的，其唯一前途，只有和全世界被压迫的劳动者、被压迫的落后民族联合在一起，推翻一切帝国主义，以分工互助的国际社会主义新世界，代替商品买卖的资本主义旧世界，民族问题便自然解决了。

对于中国文化，有人把中国史上民族的团结、印刷术与火药的发明，也排除在文化以外，把文化缩小在文学圈子里，这种误解文化的结果，在此次抗日战争中，产生了两件分不清楚的事：一是把口号喊得震天响，称为"文化人"的，却无实际行动；二是把日本侵略中国说成是"中日文化交流"，这是荒谬绝伦的。中国文化虽然有其保守倾向，每使自己民族的文化由停滞而走向衰落。但中国文化也有其优秀的传统，如重视道德修养、强调人与自然的和谐等。我们应该批判的继承中国传统文化，同时吸收外来文化的精华，以促进中国文化的发展。

张之洞的"中学为体，西学为用"，害了我们几十年。世界上没有永恒不变的文化，我们不要高唱"本位文化""东方文化"，应该以开放的心态，吸收世界各国的优秀文化，为我所用。

殖民地、半殖民地国家的民族解放运动，必须与世界范围内的反法西斯斗争和社会主义革命结合起来，才能取得胜利。南京的傀儡政府不久都会垮台！或者有人认为此次大战是轴心国和非轴心国两派帝国主义各自扩大其势力之斗争，与民族解放无关，弱小民族之参加是无意义的。这一见解是由于他们不明白民族解放自然不能够依靠帝国主义帮助而成功，也不是弱小民族自己力量可以解决的问题。而且"中立"这一名词，在现代战争中将不会再见了。缅甸人如果说：我们可和认识的魔鬼来往，而不能和不认识的天使来往，我们应该告诉他们：我们不知道现在世界上有天使，只知道你们所认识的魔鬼，比你们不认识的魔鬼，还要凶恶！中国如果有人说：帮助美国打日本是前门拒狼，后门进虎，我们应该告诉他：美国胜利了，我们如果能努力自强，不再妥协投降，可能摆脱半殖民地的地位；如果日本胜利，我们必然陷入亡国灭种的境地。

以上的话，或者有人认为是悲观论调，那只好让将来的事实来教训他。

\hfill 一九四二年二月十日

\chapter{再论世界大势（一九四二年四月十九日）}
\markboth{再论世界大势}{一九四二年四月十九日}

\quad\quad 有人批评我在《战后世界大势之轮廓》一文中所估计的国际形势，说将来只有帝国主义的天下，太悲观了。我以为最客观的估计只能如此，不必论其是否悲观。世界自前世纪以来，金融资本即已打破了民族界限，帝国主义的天下早已形成，现在不过是由七八个帝国主义国家合并为两个帝国主义集团而已。一天没有震撼全世界的大革命发生，这种状况仍会继续下去，甚至会发展到我所估计的程度。或者有人会说：此次大战如果胜利属于希特勒，英国便完了，苏联也或许会垮台，美洲的法西斯主义将起而代之，下次的世界大战即是德美之间的战争，将不是民主与纳粹之斗争，而是两派法西斯集团之间的斗争；如此则正如斯大林所说，民主自由将消失数百年才能恢复；如此则人类文化史将走入如下的道路：

\begin{center}
\small
\begin{tabularx}{\textwidth}{|l|X|}
\hline
\textbf{时代} & \textbf{政治制度} \\
\hline
上古世界 & 氏族社会民主制 \\
\hline
古代世界（希腊罗马） & 土地主、大巫、军事首领的专制 $\rightarrow$ 城市市民的民主制 \\
\hline
近代世界 & 封建诸侯及其末期的君主专制 $\rightarrow$ 资产阶级民主制 \\
\hline
未来世界 & 法西斯蒂专制 $\rightarrow$ 无产阶级民主制以至全民民主制 \\
\hline
\end{tabularx}
\end{center}

按照这一发展趋势，法西斯蒂专制将和以前的专制制度一样，普遍的发展，而且在形式上会持续一个时期，这也就是每个时代民主制向前发展之前，都要经历的一个黑暗时期；如果人们躺在幻想和安乐的温床上，听任法西斯存在发展，我们没有理由否认这一黑暗的时期到来之可能。

\hfill 一九四二年四月十九日

\chapter{被迫民族之前途（一九四二年五月十三日）}
\markboth{被迫民族之前途}{一九四二年五月十三日}

\quad\quad 被迫民族是资本主义帝国主义的产物，被迫的劳动者为其生产商品，被迫的落后民族为其提供商品市场和生产原料，这是资本主义帝国主义存在的两个支柱。

被迫民族反抗资本主义帝国主义的压迫，以至取得民族独立，那是天经地义，无可非议的。这种为民族自由而战的运动，无论为谁人所领导，民族中一切进步分子都应支持；因为不但资产阶级所领导的民族解放运动，有打击资本主义帝国主义的进步意义，即是封建王公所领导的民族解放运动，也有一定的进步作用。

但是这一运动若限制在民族斗争的范围以内，其前途如何呢？

第一：自国内方面来看，民族斗争往往会加剧国内的阶级矛盾和民族矛盾，不但不能减少，而且是增加。政治思想上的分歧，经济上的困难，如长期战争造成的封锁及通货膨胀，因为没有民主制度的保障，政治上的腐败，很容易导致汉奸和地主阶级趁机扩大势力，掠夺民财。民众一时浴血奋战，却不能用和平的手段改变这种现象，便有人大喊：这是超出了民族斗争的范围，破坏了一致对外的民族统一战线；而且在实际上也是超出了民族斗争的范围，然而如果不解决这些问题，便是民族解放的致命伤；而这些问题的解决，又必然会涉及到阶级斗争，这是不可避免的。

第二：自国际方面来看，在各派帝国主义瓜分殖民地及市场垄断化的今天，甘地一流的民族独立运动，不能依靠强国的帮助而得到自由，这还是百分之百的真理，可是没有别的强国帮助，也不能摆脱当前的强国之压迫；而且有些强国不允许你不依赖他们，这也是百分之百的真理。于是尼赫鲁等人便陷入了困境，他和甘地略有不同，即不主张依赖英国的帮助，但他也没有找到其他的出路。美国势力如果进入印度，我们知道她对于殖民地的统治，不但不比日本好，且不比英国好，菲律宾便是一个例子；然而菲律宾也不算是一个真正独立的民族国家。如果印度人以民族独立的理由（这理由当然十分正当），排斥英国势力，结果却换来了一个新主人日本，那便更糟了。甘地和尼赫鲁无论如何宣传印度受外人统治的时代已经过去了；然而他们内心也未必相信依靠自己的力量能够同时对抗英国和日本；结果不过是照旧屈服在新的主人的统治之下，执行不合作而已。这又怎样办呢？

所以我认为在资本主义帝国主义的世界，任何弱小的民族，若起来反抗，只依靠自己一个民族的力量，排除一切帝国主义之干涉，以实现孤立的民族独立，是没有出路的，其唯一前途，只有和全世界被压迫的劳动者、被压迫的落后民族联合在一起，推翻一切帝国主义，以分工互助的国际社会主义新世界，代替商品买卖的资本主义旧世界，民族问题便自然解决了。

对于我这一见解，或有人提出两个疑问：一是，落后民族如何能够联合起来？又如何能够和别的国家的劳动者及其他的弱小民族联合在一起？二是，社会主义是否包含民族解放问题？

提出第一个疑问的人，是被压迫的民族闭起了眼睛，看不出将来历史发展的新趋势。民族自己的力量，当然不足以对抗强大的帝国主义，即资本主义国家内部的无产阶级，也不能单独推翻资产阶级的统治。在今天，任何民族无论要发展资本主义，都非先取得国家独立不可，只要不是民族沙文主义者，便能够明白这一点；近百年来，资本主义帝国主义的殖民政策，已打破了各落后民族的孤立状态，此后各派帝国主义的统治形式，将由殖民政策，发展为更集中的更有组织性的集团统治，所谓《大西洋宪章》《太平洋宪章》等等，便是这种集团统治之开始。如果有一个领导集团的主要国家，如战败的德国，先和其他国家融合成一个统一的联邦。即在资本主义的国际集团内，落后民族将被吸引或被迫和其他国家合作，即这种不平等的合作也能促进集团内的各落后民族和领导国的劳动人民相互联合，这便是资本主义帝国主义自己造成了推翻自己的强大力量，没有任何民族主义的英雄能够阻止这一集团化的新趋势；而且被压迫的民族，也只有加入这种集团，才有前途。

提出第二个疑问的人，是被第二国际的机会主义理论所迷惑了。第二国际只注重在资本主义制度之下从事改良运动，所以不会计及被迫的民族解放问题，因为民族解放运动是对资本主义帝国主义的根本打击。而社会主义运动，是要根本推翻资本主义帝国主义的，所以自第一国际以来，"全世界被压迫的劳动人民"和"全世界被压迫的民族"联合起来，便是社会主义运动的重要口号。

中国的抗日战争，是世界反法西斯战争的重要组成部分，也是中国人民争取民族独立和民主自由的正义战争。我们不能把中国的抗日战争看作是单纯的民族斗争，而应该把它和世界范围内的反法西斯斗争及社会主义革命结合起来。如果我们仍然固守多年来的狭隘民族主义立场，抗日战争一开始，便以全力反对国民党政府，而不与其他抗日力量合作，中国的抗日战争便不会取得胜利。假使苏联仍然坚持以前的国际主义立场，而不是以民族利益为中心，日本便没有那么容易占领上海和南京；至少在战争初期，苏联能够给予中国更多的援助，武汉也决不致陷落，中苏一直共同抗战到最后，日本便无力进攻太平洋，蹂躏菲律宾、爪哇、缅甸等弱小民族了！当纳粹德国进攻波兰时，苏联如果仍然站在国际社会主义的立场，便不会和希特勒签订互不侵犯条约，便不会把近代民主主义抛在一边，更不会和法西斯蒂瓜分波兰！这样英法苏三国同盟便不会崩溃，希特勒便不能在西欧战场上同时得到胜利，孤立的波兰虽然失败，但其抵抗精神也会鼓舞其他国家，希特勒便没有力量征服挪威、丹麦、比利时等国。

从俄国前后立场不同其结果也不同这一连串历史事实，已是够明白国际社会主义和被压迫民族的关系了。

俄国在欧洲，毕竟也是一个比较强大的民族，她的民族政策之如何呢？她为了俄国的安全，采取了向法西斯妥协的政策，以代替反对法西斯的进攻；其结果是俄国的战争不是开始于希特勒在欧洲孤立之时，而开始于希特勒征服了欧洲各国之后，俄国不但在向法西斯妥协的时代割让给德国的半个波兰和波罗的海三国仍为希特勒所有，俄国的大部分领土和人民都陷入希特勒之手；没有英美的援助，莫斯科也未必能够守到今天。她为了俄国的安全，始终和日本保持中立，中国共产党也因此被人加以"游而不击"的恶名；其结果，明天日本仍可能协同希特勒进攻俄国，俄国民族仍处于很不安全的地位，那将不能得到中国有力的援助，因为中国被日本削弱了。所以任何落后民族，以民族政策自限，必至走上孤立无援的道路（民族政策实际上就是孤立政策），没有前途，就是俄国也不能例外。

\hfill 一九四二年五月十三日

\chapter{给Y的信（一九四二年五月十三日）}
\markboth{给Y的信}{一九四二年五月十三日}

\quad\quad Y兄：

弟陈独秀是一九四二年五月二十七日死的，这些论文和上面的意见，正如他自己所说"只是我一个人的意见，不代表任何人"（给西流等的信），而很多人不能理解和同意。

友人抄写不易，留去与否他们也不会理解和同意。第三篇文章（《再论世界大势》）我已无抄稿，望兄将原稿寄下。

诸君可以看看，由兄决定是否抄写一份给他们看。抄寄他人，可以不必，因为他们也不会理解。

是否需要抄一份给XX、XX，由他们自便。XX给你的《大公报》，请寄给我，弟不愿文章在报上刊登，而XXX已快要离开了。XX已赴印度，前已去信，想必已经收到了。兹寄上一文（即《被迫民族之前途》），算是前三篇（按：《我的根本意见》《战后世界大势之轮廓》《再论世界大势》）的结论，更详细的内容，以后再谈。

顺祝健康！

\hfill 弟独秀\quad 手启\quad 五月十三日（一九四二年）

（编者按：此信为陈独秀晚年最后几封信之一，反映了他晚年的孤独和坚持自己见解的决心。）

% ==================== 文档结束 ====================
\end{document}
